\documentclass[11pt]{article}
\usepackage{amsmath,hyperref}
\title{The Last Useful Notice for a Moving Inspection Window}
\date{2026-11-27}
\begin{document}
\maketitle

\paragraph{Audit.}
Two model robots begin at \(0,4\). The first stays still and the second
has speed at most one, order and separation in \([1,6]\), and the
modeled radio link given by the upper bound. Each footprint has sensing
radius two; ordinary single coverage must hold at every point of the
owed interval, of width one. Only the horizon \([0,2]\) is required,
not the return of the previous example. The locally readable inputs
are the robot's own state, the common time and a truthful qualified
notice at a prescribed cut \(r\). There is no earlier evidence
distinguishing the requirements. A future mission designation is
not inferred from local geometry. The proof models the information
cut; the fallible physical notice channel, clock error, and who can
produce the notice are outside this specialization.

\paragraph{Two individually specified worlds.}
The far-world window starts at \(s_F(t)=h t/2\), \(h\ge0\). The other
starts at \(s_L(t)=\min(3t/2,3/2)\). In the second world the two
footprints can stay in their initial configuration. For \(h>2\)
the previous seam argument requires the movable robot to be at
most \(4\) at \(a=4/h\), when the far interval has just finished
crossing position \(2\). It can achieve the terminal requirement
exactly when \(h+4/h\le7\) by holding first then advancing at unit
speed (for \(h\le2\) static positions suffice). Thus the physical
limit is \(H=(7+\sqrt{33})/2\). With \(h\le5\) both worlds can use
static positions without notice. The nontrivial common-prefix
question is \(5<h\le H\). The policies are deterministic and
must guarantee each realized world; private randomness cannot
certify a randomized failure branch for this hard service law.

\paragraph{Sharp reveal-time bound.}
In the late-window world, the interval straddles position \(2\)
from just after \(c=2/3\) through the deadline. Thus the second
robot must be at most \(4\) from \(c\), including \(c\) by continuity.
At an earlier reveal it can be ahead of \(4\) only by at most
its remaining retreat distance \(c-r\). In the far world, reaching
the terminal mobile position at least \(h-1\) by time \(2\)
requires its cut position to be at least \(h-3+r\). The same
prefix produces the same cut position. Consequently
\[
  h-3+r\;\le\;4+\max(0,c-r).
\]
Solving gives the maximum admissible reveal time
\[
 r_*(h)=
 \begin{cases}
 7-h,&5<h\le19/3,\\
 (23/3-h)/2,&19/3<h\le H.
 \end{cases}
\]
This obstruction holds over all continuous causal speed-feasible
prefixes, not just controllers that wait at their initial positions.

For sufficiency, in the first case hold at \(4\) until notice.
The far branch then advances at unit speed starting at
\(\max(r,a)\); the late branch holds. Both \(r\) and \(a\) are no
later than \(7-h\), by the cutoff and physical feasibility.
The early window crossing remains covered by the touching initial
footprints; after crossing, the far left endpoint advances faster
than the moving footprint's left edge. The tight right-hand
endpoint condition is at the deadline. These arguments also
cover \(r\le7-h\) in the second case. If \(7-h<r\le r_*\),
start advancing at \(7-h\) before the notice, which is no earlier
than \(a\). The shared prefix remains covered in the late world
because \(r<c\). At notice the far branch continues at unit speed;
the late branch retreats at unit speed to \(4\). The retreat can
finish by \(c\) precisely because of the upper inequality at the
cut. The gap stays within \([1,6]\). For example at \(h=6.37\)
the best wait-until-notice variant allows \(63/100\), and the
prefix with a contingent retreat allows \(389/600\), an extra
\(11/600\). At \(h=6.38\) a failure is instead physical, even
with full advance information. Neither delay number is a
measured protocol budget.

\paragraph{Exact reference and decoder.}
On each side of the switching times \(s=1,2\), the allowed mobile
positions are one closed interval given by affine bounds after
taking intersections with the gaps. At shared knots retain both
one-sided constraints by continuity. In the open strata these
intervals express actual coverage by the union of the two
ordered footprints; affine interpolation of seam positions
satisfies both bounds throughout a stratum.
A forward reachable interval expands by the segment duration
and intersects the next admissible interval. Backward propagation
from each required completion yields two viable seam intervals.
The common prefix uses the intersection of the two worlds'
requirements and must land in both suffix intervals.
Reconstruct earlier and later seam positions and use affine
segments; the first post-notice point is the same physical point
in both branches.

The direct comparator assembles the complete time-knot network
with two separate suffixes and a shared seam and checks its
difference constraints by exact negative-cycle detection. The
answer agrees on 3571 \((h,r)\) probes, plus 170 individual-world
checks. There are 2535 constructive positive shared policies,
checked against 94620 exact footprint-order event or open-stratum
samples by an independent continuous-coverage criterion.
66 rational cutoff equalities are tested. Lean checks exact
interval expansion and scalar cut inequalities; it is not a
mechanized verification of optical assumptions or the runtime.
The constructive decoder is an instance of exact planning with
nonanticipativity, not a claimed improvement over a
same-information direct solver. Motion clearance and timely
mission authority remain separate evidence obligations.
\end{document}
